\documentclass[11pt,a4paper]{article}
\usepackage{../common/td_style}

\begin{document}

\makeuniversityheader{Corrigé des Travaux Dirigés : Série 00}{Introduction aux Biostatistiques \& Rappels Fondamentaux}{\textcolor{BioTeal}{\textbf{CORRIGÉ DÉTAILLÉ}}}

\begin{solutionbox}{Exercice 1 : Typologie des Variables en Biochimie Clinique}
\begin{enumerate}[label=\textbf{\alph*.}]
  \item \textbf{Nature exacte des variables :}
  \begin{itemize}
    \item \textbf{Sexe du patient :} Variable \textbf{qualitative nominale} binaire (deux modalités non hiérarchisées : Masculin / Féminin).
    \item \textbf{Triglycérides (g/L) :} Variable \textbf{quantitative continue} (résultat de mesure biochimique sur une échelle d'intervalles/ratios réels).
    \item \textbf{Stade de fibrose hépatique (F0 à F4) :} Variable \textbf{qualitative ordinale} (il existe un ordre biologique strict de sévérité, mais la distance entre F0 et F1 n'est pas nécessairement égale à celle entre F3 et F4).
    \item \textbf{Nombre d'épisodes d'ictère :} Variable \textbf{quantitative discrète} (valeurs entières $\in \mathbb{N}$ issues d'un décompte).
    \item \textbf{Absorbance à $540\,\text{nm}$ :} Variable \textbf{quantitative continue} (mesure optique continue, échelle continue positive).
    \item \textbf{Polymorphisme \textit{PNPLA3} :} Variable \textbf{qualitative nominale} (ou ordinale si codée par le nombre d'allèles mutés 0, 1, 2).
  \end{itemize}

  \item \textbf{Indicateurs statistiques recommandés :}
  \begin{itemize}
    \item \textit{Variables qualitatives (Sexe, Stades, Génotypes) :} Mode et pourcentages (fréquences relatives). Pour le stade ordinale, la \textbf{Médiane} et les rangs sont appropriés.
    \item \textit{Variables quantitatives continues symétriques/normales :} \textbf{Moyenne $\pm$ Écart-type ($SD$)}.
    \item \textit{Variables quantitatives asymétriques ou avec outliers (ex: épisodes d'ictère, triglycérides) :} \textbf{Médiane [Intervalle Interquartile $Q_1 - Q_3$]}.
  \end{itemize}
\end{enumerate}
\end{solutionbox}

\begin{solutionbox}{Exercice 2 : Variabilité Expérimentale \& Mesures de Dispersion}
\begin{enumerate}[label=\textbf{\alph*.}]
  \item \textbf{Calculs pas-à-pas ($N = 6$) :}
  \begin{itemize}
    \item Somme des valeurs : $\sum x_i = 38.2 + 41.5 + 39.8 + 42.0 + 37.5 + 41.0 = 240.0\,\text{g/L}$.
    \item \textbf{Moyenne empirique :} $\bar{x} = \frac{240.0}{6} = \mathbf{40.00\,\text{g/L}}$.
    \item Écarts à la moyenne $(x_i - \bar{x})$ : $-1.8, \, +1.5, \, -0.2, \, +2.0, \, -2.5, \, +1.0$.
    \item Carrés des écarts $(x_i - \bar{x})^2$ : $3.24 + 2.25 + 0.04 + 4.00 + 6.25 + 1.00 = 16.78$.
    \item \textbf{Variance sans biais ($s^2$) :} $s^2 = \frac{\sum (x_i - \bar{x})^2}{N - 1} = \frac{16.78}{5} = \mathbf{3.356\,(\text{g/L})^2}$.
    \item \textbf{Écart-type ($s$) :} $s = \sqrt{3.356} = \mathbf{1.832\,\text{g/L}}$.
  \end{itemize}

  \item \textbf{Coefficient de variation ($CV$) :}
  \begin{equation*}
    CV = \frac{s}{\bar{x}} \times 100 = \frac{1.832}{40.00} \times 100 = \mathbf{4.58\%}
  \end{equation*}
  Puisque $CV = 4.58\% < 5\%$, la méthode satisfait le critère usuel de répétabilité en biochimie clinique.

  \item \textbf{Différence $N$ vs $N-1$ (Correction de Bessel) :}
  La division par $N$ sous-estime systématiquement la variance de la population car les écarts sont calculés par rapport à la moyenne de l'échantillon $\bar{x}$ (qui minimise mathématiquement la somme des carrés) et non par rapport à la vraie moyenne inconnue $\mu$. Diviser par $N-1$ (les degrés de liberté restants) compense exactement ce biais.
\end{enumerate}
\end{solutionbox}

\begin{solutionbox}{Exercice 3 : SD versus SEM \& Intervalle de Confiance à 95\%}
\begin{enumerate}[label=\textbf{\alph*.}]
  \item \textbf{Erreur standard de la moyenne ($SEM$) :}
  \begin{equation*}
    SEM = \frac{s}{\sqrt{N}} = \frac{20.0}{\sqrt{25}} = \frac{20.0}{5} = \mathbf{4.00\,\text{U/mg}}
  \end{equation*}

  \item \textbf{Différence biologique fondamentale :}
  \begin{itemize}
    \item \textbf{Écart-type ($SD = 20.0$) :} Mesure la \textbf{dispersion biologique réelle} des individus autour de la moyenne. Il ne diminue pas lorsque la taille d'échantillon augmente.
    \item \textbf{Erreur standard ($SEM = 4.0$) :} Mesure la \textbf{précision de l'estimation de la moyenne}. C'est une métrique d'inférence. L'utiliser abusivement pour faire paraître les barres d'erreur plus petites est une pratique trompeuse très fréquente en biologie.
  \end{itemize}

  \item \textbf{Intervalle de Confiance à 95\% :}
  \begin{align*}
    \text{IC}_{95\%} &= \bar{x} \pm t_{0.05, \, 24} \times SEM \\
    &= 142.0 \pm (2.064 \times 4.00) = 142.0 \pm 8.26 = \mathbf{[133.74 \,;\, 150.26]\,\text{U/mg}}
  \end{align*}
  \textit{Interprétation :} Si nous répétions cette expérience 100 fois dans les mêmes conditions, 95\% des intervalles calculés contiendraient la vraie moyenne biologique $\mu$ de la population.

  \item \textbf{Effet de $N = 100$ :}
  $SEM_{100} = \frac{20}{\sqrt{100}} = 2.0\,\text{U/mg}$ (divisé par 2). L'intervalle de confiance devient deux fois plus étroit, illustrant le gain de précision lié à la taille d'échantillon.
\end{enumerate}
\end{solutionbox}

\begin{solutionbox}{Exercice 4 : Logique du Test d'Hypothèse \& Démystification de la $p$-value}
\begin{enumerate}[label=\textbf{\alph*.}]
  \item \textbf{Formulation des hypothèses :}
  \begin{itemize}
    \item $H_0$ (Hypothèse nulle) : $\mu_1 = \mu_2$ (l'extrait végétal n'a aucun effet hypoglycémiant ; la différence constatée est purement due aux fluctuations d'échantillonnage).
    \item $H_1$ (Hypothèse alternative bilatérale) : $\mu_1 \neq \mu_2$ (l'extrait modifie significativement la glycémie).
  \end{itemize}

  \item \textbf{Définition probabiliste exacte de $p = 0.0042$ :}
  C'est la probabilité d'observer une différence de moyenne au moins aussi grande que celle constatée ($|3.20 - 2.65| = 0.55\,\text{g/L}$), \textbf{sous l'hypothèse stricte que le traitement n'a aucun effet ($H_0$ vraie)}.
  $p = P(\text{Données observées ou plus extrêmes} \mid H_0 \text{ est vraie}) = 0.0042$.

  \item \textbf{Pourquoi l'affirmation 99.58\% est fausse :}
  La $p$-value n'est \textbf{PAS} la probabilité que $H_0$ ou $H_1$ soit vraie ($P(H_0 \mid \text{Données}) \neq P(\text{Données} \mid H_0)$). Confondre les deux est l'erreur du procureur (inversion conditionnelle). Pour calculer la probabilité d'efficacité, il faudrait une approche bayésienne avec une probabilité a priori.

  \item \textbf{Conclusion biologique :}
  Puisque $p = 0.0042 < \alpha = 0.05$, on \textbf{rejette l'hypothèse nulle $H_0$}. L'extrait végétal induit une diminution statistiquement significative de la glycémie ($p < 0.01$).
\end{enumerate}
\end{solutionbox}

\begin{solutionbox}{Exercice 5 : Contrôle de Normalité \& Détection d'Outliers}
\begin{enumerate}[label=\textbf{\alph*.}]
  \item \textbf{Calcul des quartiles sur les données ordonnées ($N=8$) :}
  Données triées : $21.8, \quad 22.1, \quad 22.5, \quad 22.9, \quad 23.0, \quad 23.4, \quad 24.1, \quad 48.7$.
  \begin{itemize}
    \item \textbf{Médiane} (moyenne des 4\textsuperscript{e} et 5\textsuperscript{e} valeurs) : $\frac{22.9 + 23.0}{2} = \mathbf{22.95\%}$.
    \item \textbf{Premier quartile $Q_1$} (médiane de la 1\textsuperscript{re} moitié) : $\frac{22.1 + 22.5}{2} = \mathbf{22.30\%}$.
    \item \textbf{Troisième quartile $Q_3$} (médiane de la 2\textsuperscript{de} moitié) : $\frac{23.4 + 24.1}{2} = \mathbf{23.75\%}$.
    \item \textbf{Écart interquartile ($IQR$)} : $IQR = Q_3 - Q_1 = 23.75 - 22.30 = \mathbf{1.45\%}$.
  \end{itemize}

  \item \textbf{Règle de Tukey pour les barrières d'outliers :}
  \begin{itemize}
    \item Barrière inférieure : $Q_1 - 1.5 \times IQR = 22.30 - (1.5 \times 1.45) = 22.30 - 2.175 = \mathbf{20.125\%}$.
    \item Barrière supérieure : $Q_3 + 1.5 \times IQR = 23.75 + 2.175 = \mathbf{25.925\%}$.
  \end{itemize}
  La valeur $48.7\%$ est très largement supérieure à $25.925\%$ : \textbf{c'est formellement une valeur aberrante extrême (outlier)}.

  \item \textbf{Attitude scientifique du biochimiste :}
  Ne jamais supprimer une valeur sans traçabilité !
  \begin{enumerate}
    \item Vérifier le cahier de laboratoire (erreur de pipetage, bulle d'air dans la cuve spectrophotométrique, turbidité anormale).
    \item Si une erreur technique avérée est identifiée, la valeur peut être écartée en documentant la raison.
    \item Si aucune erreur technique n'est trouvée, présenter les analyses avec et sans cette observation (analyse de sensibilité).
  \end{enumerate}

  \item \textbf{Test non-paramétrique alternatif :}
  En cas de non-normalité ou d'outlier non éliminable, le \textbf{test $U$ de Mann-Whitney} (ou test des rangs de Wilcoxon pour séries indépendantes) doit être utilisé à la place du test $t$ de Student.
\end{enumerate}
\end{solutionbox}

\begin{solutionbox}{Exercice 6 : Théorème de Bayes \& Performances Diagnostiques (ELISA)}
\begin{enumerate}[label=\textbf{\alph*.}]
  \item \textbf{Probabilité de faux positif :}
  $P(+ \mid \bar{M}) = 1 - Sp = 1 - 0.95 = \mathbf{0.05 \quad (5\%)}$.

  \item \textbf{Probabilité totale d'un résultat positif $P(+)$ :}
  D'après la formule des probabilités totales :
  \begin{align*}
    P(+) &= P(+ \mid M) \cdot P(M) + P(+ \mid \bar{M}) \cdot P(\bar{M}) \\
    &= (0.98 \times 0.01) + (0.05 \times 0.99) = 0.0098 + 0.0495 = \mathbf{0.0593 \quad (5.93\%)}
  \end{align*}

  \item \textbf{Valeur Prédictive Positive ($VPP$) par le Théorème de Bayes :}
  \begin{equation*}
    VPP = P(M \mid +) = \frac{P(+ \mid M) \cdot P(M)}{P(+)} = \frac{0.0098}{0.0593} = \mathbf{0.1653 \quad (16.53\%)}
  \end{equation*}

  \item \textbf{Explication du paradoxe des faux positifs :}
  Bien que le test soit excellent ($Se = 98\%$, $Sp = 95\%$), la maladie est rare ($1\%$). Sur un échantillon de $10\,000$ personnes :
  \begin{itemize}
    \item $100$ malades $\implies 100 \times 0.98 = \mathbf{98}$ vrais positifs ($2$ faux négatifs).
    \item $9\,900$ personnes saines $\implies 9\,900 \times 0.05 = \mathbf{495}$ faux positifs !
  \end{itemize}
  Le nombre de faux positifs ($495$) est cinq fois supérieur au nombre de vrais positifs ($98$). Ainsi, parmi les $593$ individus testés positifs, seuls $16.5\%$ sont réellement infectés. \textbf{Conséquence clinique :} En dépistage de masse, tout test ELISA positif doit impérativement faire l'objet d'un test de confirmation plus spécifique (Western Blot, RT-PCR).

  \item \textbf{Valeur Prédictive Négative ($VPN$) :}
  \begin{align*}
    VPN &= P(\bar{M} \mid -) = \frac{Sp \cdot (1 - P(M))}{Sp \cdot (1 - P(M)) + (1 - Se) \cdot P(M)} \\
    &= \frac{0.95 \times 0.99}{(0.95 \times 0.99) + (0.02 \times 0.01)} = \frac{0.9405}{0.9407} = \mathbf{99.98\%}
  \end{align*}
  Un résultat négatif est d'une fiabilité quasi-absolue pour écarter l'infection.

  \item \textbf{Effet de la prévalence hospitalière ($P(M) = 30\%$) :}
  \begin{align*}
    P(+) &= (0.98 \times 0.30) + (0.05 \times 0.70) = 0.2940 + 0.0350 = 0.3290 \\
    VPP &= \frac{0.2940}{0.3290} = \mathbf{89.36\%}
  \end{align*}
  \textit{Conclusion fondamentale :} La sensibilité et la spécificité sont des propriétés intrinsèques du réactif biologique, mais la $VPP$ dépend directement de la prévalence épidémiologique de la population explorée.
\end{enumerate}
\end{solutionbox}

\begin{solutionbox}{Exercice 7 : Lois de Probabilités en Microbiologie \& Génétique}
\textbf{Partie 1 : Contrôle de Stérilité (Loi de Poisson $\lambda = 3.0$)}
\begin{enumerate}[label=\textbf{\alph*.}]
  \item \textbf{Justification biologique :} Le dénombrement d'UFC représente des événements rares, discrets ($X \in \mathbb{N}$), indépendants, survenant de manière aléatoire et uniforme dans un volume continu ($100\,\mu\text{L}$).
  \item \textbf{Boîte parfaitement stérile ($X = 0$) :}
  $P(X = 0) = \frac{e^{-3.0} \cdot 3.0^0}{0!} = e^{-3.0} = \mathbf{0.0498 \quad (4.98\%)}$.
  \item \textbf{Exactement 2 colonies ($X = 2$) :}
  $P(X = 2) = \frac{e^{-3.0} \cdot 3.0^2}{2!} = \frac{9 \cdot 0.04979}{2} = \mathbf{0.2240 \quad (22.40\%)}$.
  \item \textbf{Contamination notable ($X \ge 4$) :}
  \begin{align*}
    P(X \ge 4) &= 1 - [P(0) + P(1) + P(2) + P(3)] \\
    P(1) &= \frac{e^{-3} \cdot 3^1}{1!} = 0.1494, \quad P(3) = \frac{e^{-3} \cdot 27}{6} = 0.2240 \\
    P(X \ge 4) &= 1 - [0.0498 + 0.1494 + 0.2240 + 0.2240] = 1 - 0.6472 = \mathbf{0.3528 \quad (35.28\%)}
  \end{align*}
\end{enumerate}

\textbf{Partie 2 : Sélection de Mutants (Loi Binomiale $n=10, \, p=0.05$)}
\begin{enumerate}[label=\textbf{\alph*.}]
  \setcounter{enumi}{4}
  \item \textbf{Justification \& Paramètres :} Succession de $n = 10$ tirages indépendants de Bernoulli avec probabilité constante de succès $p = 0.05$.
  $\mathbb{E}(Y) = n \cdot p = 10 \times 0.05 = \mathbf{0.50 \text{ mutant}}$, et $\sigma_Y = \sqrt{n \cdot p \cdot (1-p)} = \sqrt{10 \times 0.05 \times 0.95} = \sqrt{0.475} = \mathbf{0.689}$.
  \item \textbf{Exactement un clone résistant ($Y = 1$) :}
  $P(Y = 1) = \binom{10}{1} \cdot (0.05)^1 \cdot (0.95)^9 = 10 \times 0.05 \times 0.6302 = \mathbf{0.3151 \quad (31.51\%)}$.
  \item \textbf{Au moins un clone résistant ($Y \ge 1$) \& Taille d'échantillon :}
  $P(Y \ge 1) = 1 - P(Y = 0) = 1 - (0.95)^{10} = 1 - 0.5987 = \mathbf{0.4013 \quad (40.13\%)}$.
  Pour obtenir au moins un mutant avec une confiance $\ge 95\%$ :
  \begin{align*}
    1 - (0.95)^n &\ge 0.95 \implies (0.95)^n \le 0.05 \\
    &\implies n \ge \frac{\ln(0.05)}{\ln(0.95)} = \frac{-2.9957}{-0.05129} \approx 58.4 \implies \mathbf{n = 59 \text{ clones}}
  \end{align*}
\end{enumerate}
\end{solutionbox}

\begin{solutionbox}{Exercice 8 : Loi Normale, $Z$-Score \& Intervalles de Référence}
\begin{enumerate}[label=\textbf{\alph*.}]
  \item \textbf{Intervalle de référence physiologique à 95\% :}
  $[\mu \pm 1.96\sigma] = [0.95 - (1.96 \times 0.10) \,;\, 0.95 + (1.96 \times 0.10)] = [0.95 - 0.196 \,;\, 0.95 + 0.196] = \mathbf{[0.754 \,;\, 1.146]\,\text{g/L}}$.
  95\% de la population saine présente une glycémie comprise entre $0.75$ et $1.15\,\text{g/L}$.

  \item \textbf{Calcul du $Z$-score du patient ($X = 1.22\,\text{g/L}$) :}
  \begin{equation*}
    Z = \frac{X - \mu}{\sigma} = \frac{1.22 - 0.95}{0.10} = \frac{0.27}{0.10} = \mathbf{+2.70}
  \end{equation*}
  La glycémie de ce patient est située à $2.70$ écarts-types au-dessus de la moyenne de référence.

  \item \textbf{Proportion de sujets sains avec $X \ge 1.22\,\text{g/L}$ :}
  $P(X \ge 1.22) = P(Z \ge 2.70) = 1 - \Phi(2.70) = 1 - 0.9965 = \mathbf{0.0035 \quad (0.35\%)}$.
  Seulement $3.5$ personnes saines sur $1\,000$ atteignent un tel niveau. Cette valeur est nettement pathologique et justifie le diagnostic d'hyperglycémie.

  \item \textbf{Asymétrie des biomarqueurs \& Transformation logarithmique :}
  Pour des enzymes comme l'ALAT, la présence de queues de distribution vers les valeurs très élevées rend la distribution fortement asymétrique (non normale). Appliquer $\mu \pm 2\sigma$ aboutirait à une borne inférieure absurde (parfois négative) et tronquerait les valeurs réelles. On applique en routine une \textbf{transformation logarithmique} $Y = \log_{10}(X)$ pour symétriser les données et rétablir la normalité gaussienne.
\end{enumerate}
\end{solutionbox}

\begin{solutionbox}{Exercice 9 : Puissance Statistique, Risque $\beta$ \& Éthique 3R}
\begin{enumerate}[label=\textbf{\alph*.}]
  \item \textbf{Risques $\alpha$ et $\beta$ :}
  \begin{itemize}
    \item \textbf{Risque $\alpha$ (Première espèce) :} Probabilité de rejeter $H_0$ à tort (Faux positif, conclure à un effet qui n'existe pas).
    \item \textbf{Risque $\beta$ (Seconde espèce) :} Probabilité de conserver $H_0$ à tort (Faux négatif, manquer un effet thérapeutique réel).
  \end{itemize}

  \item \textbf{Puissance statistique ($1 - \beta$) :}
  C'est la probabilité de détecter un effet réel d'une taille donnée s'il existe. Le consensus scientifique international (ex: FDA, ARRIVE guidelines) impose une puissance minimale de $\mathbf{80\%}$ ($1 - \beta \ge 0.80$), et idéalement $90\%$ pour les essais pivots.

  \item \textbf{Conséquences d'une puissance de 25\% ($N = 3$) :}
  \begin{itemize}
    \item Le risque de faux négatif est de $\beta = 1 - 0.25 = \mathbf{75\%}$. Le chercheur a 3 chances sur 4 de ne rien observer et d'abandonner une molécule pourtant active.
    \item \textbf{Violation éthique des 3R :} Le principe de « Réduction » ne signifie pas réduire le nombre d'animaux au point de rendre l'expérience inutile. Une étude sous-puissancée entraîne le sacrifice pur et simple d'animaux sans possibilité de tirer une conclusion scientifique valide.
  \end{itemize}

  \item \textbf{Trois leviers pour augmenter la puissance sans gonfler le nombre d'animaux :}
  \begin{enumerate}
    \item \textbf{Réduire la variance d'erreur résiduelle ($\sigma^2$) :} Standardisation stricte de l'animalerie (cycle lumière, température, lignées pures), protocoles de pipetage automatisés.
    \item \textbf{Augmenter l'amplitude de l'effet ($\Delta$) :} Optimiser la dose, le temps d'exposition ou la sensibilité du test biologique mesuré.
    \item \textbf{Utiliser un plan expérimental apparié ou à mesures répétées :} Chaque animal servant de son propre témoin (mesures avant et après traitement), la variabilité inter-individuelle est mathématiquement éliminée de l'erreur résiduelle.
  \end{enumerate}
\end{enumerate}
\end{solutionbox}

\end{document}
